%============================================================================
% prelude.tex
%   - titlepage
%   - table of contents, list of tables and list of figures
%   - abstract
%============================================================================

\clearpage%
\pagenumbering{roman}  % This makes the page numbers Roman (i, ii, etc)
%--------------------------------------------------------------------%
% TITLE PAGE
%   - define \title{} \author{} \date{}

\title{Study on Fruit Quality Inspection Based on Color Image Processing}
\author{Chitra Sures}

\date{October 2011}

\department{Department of Electronics Engineering}

%    - Set the guide's name
\setguide{Prof. M. D. Patil}
%    - Set the coguide's name (if you have one)
%%\setcoguide{Prof Amitabha Sanyal}

%   - once the above are defined, use \maketitle to generate the titlepage
\maketitle

%--------------------------------------------------------------------
\cleardoublepage \nonumber
\begin{center}
\includegraphics{raitlogo.eps}\\
Ramrao Adik Education Society's\\
\large {\textbf{Ramrao Adik Institute of Technology}}\\
\small(Affiliated to the University of Mumbai)\\
Dr. D. Y. Patil Vidyanagar,Sector 7, Nerul, Navi Mumbai 400
706.\\\vspace{0.3in} \Large CERTIFICATE
\end{center}
\vspace{0.1in}
\begin{center}
\large This to certify that Special Topic Seminar entitled\\
\textbf{\lq\lq Study on Fruit Quality Inspection Based on Color Image Processing\rq\rq} \\submitted by
\begin{center}
Hemlata Kadam (10EL2007)  \\
\end{center}
is approved for the degree of \\ \textbf {Master of Engineering}\\ in\\ \textbf {Electronics Engineering}.\\

 %\vspace{0.5in}
%-----------------------------------------------------------------
%\begin{tabular}{ccc}
%      \rule{6cm}{1sp}                &\rule{10mm}{0pt}& \rule{6cm}{1sp}
%      \\
%       Guide             &&Head of Department \\\\\
%      \rule{6cm}{1sp}                && \rule{6cm}{1sp} \\
%      Examiner 1                 && Examiner 2 \\\\\\
%      Principal: \rule{4cm}{1sp} && \\ && \\
%       &&
%    \end{tabular}
\vspace{1in}
\begin{tabular}{ccc}
      \rule{5cm}{1sp}                &\rule{10mm}{0pt}& \rule{5cm}{1sp}
      \\\vspace{0.5in}
      Examiner 1                 && Examiner 2 \\
      \rule{5cm}{1sp}                && \rule{5cm}{1sp} \\ \vspace{0.5in}
       Guide             && Project Coordinator \\
      \rule{5cm}{1sp}                && \rule{5cm}{1sp} \\
      Head of Department                && Principal \\\\
        %&\rule{5cm}{1sp}& \\
%       &Principal&
    \end{tabular}
\end{center}
%-----------------------------------------------------------------
\cleardoublepage
%--------------------------------------------------------------------%


% ABSTRACT
\begin{abstract}
%\input {abstract}
In this paper, a new method of inspecting fruit
quality is proposed based on color image processing. After an
image of fruits is taken, white balance is performed.Then the
image is transferred from the RGB color model to the HSI
color model. Its simplified histograms of hue H and saturation
S are calculated as the input of a designed BP network. The
output of the BP network is the quality description of the
inspected fruits.After training, the quality of fruits is inspected by the
BP network according to the simplified histograms of H and S
of their color image. Experiments are conducted with the
quality inspection of bananas. Experiment results show the
feasibility and reliability of proposed method.
\end{abstract}

%--------------------------------------------------------------------%
% CONTENTS, TABLES, FIGURES
\tableofcontents \listoftables \listoffigures

%--------------------------------------------------------------------%
%  Single counter for theorems and theorem-like environments:
\newtheorem{theorem}{Theorem}[chapter]
\newtheorem{assertion}[theorem]{Assertion}
\newtheorem{claim}[theorem]{Claim}
\newtheorem{conjecture}[theorem]{Conjecture}
\newtheorem{corollary}[theorem]{Corollary}
\newtheorem{definition}[theorem]{Definition}
\newtheorem{example}[theorem]{Example}
\newtheorem{figger}[theorem]{Figure}
\newtheorem{lemma}[theorem]{Lemma}
\newtheorem{prop}[theorem]{Proposition}
\newtheorem{remark}[theorem]{Remark}

%--------------------------------------------------------------------%
% Make the page numbers Arabic (1, 2, etc)
\cleardoublepage%
\pagenumbering{arabic}
